\cleardoublepage
\section*{CHƯƠNG 1. TỔNG QUAN VỀ BÀI TOÁN PHÁT HIỆN GIAN LẬN THẺ TÍN DỤNG}
\addcontentsline{toc}{section}{\numberline{}CHƯƠNG 1. TỔNG QUAN VỀ BÀI TOÁN PHÁT HIỆN GIAN LẬN THẺ TÍN DỤNG}

\setcounter{section}{1}
\setcounter{subsection}{0}
\setcounter{figure}{0}
\setcounter{table}{0}


Trong chương này, em sẽ tập trung giới thiệu về bài toán phát hiện gian lận thẻ tín dụng và các thách thức đặc thù mà bài toán này đang đặt ra. Trước tiên, em sẽ trình bày tổng quan về bản chất và tầm quan trọng của việc phát hiện gian lận trong bối cảnh thanh toán điện tử phát triển mạnh mẽ. Theo \cite{nilson2020fraud}, thiệt hại toàn cầu do gian lận thẻ tín dụng đã cán mốc \$35 tỉ vào năm 2020, cho thấy mức độ khẩn thiết của việc xây dựng hệ thống giám sát gian lận thời gian thực. Gần đây, thống kê từ Ngân hàng Trung ương châu Âu (ECB, 2024) cho thấy giao dịch gian lận dạng \textit{card-not-present} chiếm đến 80\% tổng số vụ gian lận thẻ, khẳng định xu hướng gia tăng phức tạp và tinh vi của loại hình này. Một ví dụ điển hình là vụ rò rỉ dữ liệu của hơn một triệu tài khoản ngân hàng năm 2023, dẫn đến thiệt hại hàng triệu USD và gây tổn hại nghiêm trọng đến uy tín của các tổ chức tài chính. Trong bối cảnh chuyển đổi số mạnh mẽ của ngành ngân hàng tại Việt Nam và quốc tế, các quy định pháp lý như PCI-DSS và GDPR cũng ngày càng nghiêm ngặt hơn, nhấn mạnh tầm quan trọng của các hệ thống bảo mật và phát hiện gian lận tiên tiến. Tiếp theo, chương sẽ phân tích các thách thức chính của bài toán như sự mất cân bằng lớp, tính tuần tự của dữ liệu và quan hệ phức tạp giữa các đặc trưng. Từ đó, em sẽ giới thiệu mô hình MaRENet – một phương pháp học sâu được đề xuất nhằm giải quyết các thách thức trên bằng cách kết hợp các cơ chế Transactional Expansion, Channel-wise Feature Mask, RNN và Focal Loss. Cuối cùng, chương sẽ trình bày mục tiêu và phương pháp nghiên cứu chính của đồ án.

\subsection{Giới thiệu về phát hiện gian lận thẻ tín dụng}

Phát hiện gian lận thẻ tín dụng là một nhánh nghiên cứu quan trọng trong lĩnh vực tài chính – ngân hàng, đóng vai trò thiết yếu trong việc đảm bảo an toàn cho hệ thống thanh toán điện tử. Trong bối cảnh giao dịch trực tuyến ngày càng phổ biến, các hành vi gian lận cũng trở nên tinh vi và khó phát hiện hơn. Gian lận thẻ tín dụng không chỉ gây tổn thất tài chính lớn cho các tổ chức tín dụng, mà còn ảnh hưởng nghiêm trọng đến niềm tin của người tiêu dùng vào các hệ thống thanh toán số. Do đó, việc phát triển các mô hình có khả năng phát hiện sớm và chính xác các hành vi gian lận là yêu cầu cấp thiết đối với các hệ thống tài chính hiện đại.

Khác với các bài toán học máy thông thường, dữ liệu giao dịch tài chính có nhiều đặc điểm đặc biệt như: tỷ lệ giao dịch gian lận rất nhỏ so với tổng số giao dịch (mất cân bằng lớp nghiêm trọng), giao dịch thường diễn ra theo chuỗi thời gian (tính tuần tự), và tồn tại các mối quan hệ phức tạp giữa người dùng, thời gian, thiết bị, địa điểm,... làm gia tăng độ khó của việc phát hiện. Các nghiên cứu gần đây đã chỉ ra rằng các mô hình sâu như LSTM~\cite{hochreiter1997long}, Transformer~\cite{vaswani2017attention} hay STGN~\cite{xie2023spatial} đều đạt hiệu quả cao trong việc phát hiện các mẫu gian lận tinh vi. Tuy nhiên, những phương pháp này cũng đặt ra thách thức không nhỏ về khả năng xử lý thời gian thực và chi phí tính toán lớn.


Hình \ref{fraud_overview} mô tả sơ đồ tổng thể của một hệ thống phát hiện gian lận thẻ tín dụng điển hình trong thực tế, gồm các bước từ tiền xử lý dữ liệu, huấn luyện và lưu trữ mô hình, đến dự đoán, điều tra và phản hồi để cải thiện mô hình liên tục.

\begin{figure}[!ht]
\centering
\includegraphics[width=0.9\textwidth]{Images/system_pipeline.png}
\caption[Quy trình phát hiện gian lận thẻ tín dụng.]{\textbf{Quy trình phát hiện gian lận thẻ tín dụng}~\cite{khalid2024enhancing}}
\label{fraud_overview}
\end{figure}

Cụ thể, hệ thống này vận hành qua sáu bước chính: (1) tiền xử lý và phân tách dữ liệu huấn luyện, kiểm thử; (2) huấn luyện mô hình dự đoán; (3) lưu mô hình đã huấn luyện; (4) ứng dụng mô hình để đánh giá mức độ rủi ro của giao dịch trực tuyến; (5) gửi các giao dịch đáng ngờ đến bộ phận điều tra chuyên trách; và (6) cập nhật kết quả điều tra để cải thiện hiệu suất mô hình.

Các phương pháp phát hiện gian lận truyền thống thường dựa trên luật cố định hoặc kỹ thuật thống kê đơn giản, nhưng thiếu khả năng thích ứng với các mẫu hành vi mới và linh hoạt theo thời gian. Trong khi đó, các mô hình học sâu (Deep Learning) cho thấy tiềm năng mạnh mẽ trong việc học các đặc trưng phức tạp, phát hiện các mẫu gian lận ẩn mà các phương pháp thông thường khó nắm bắt.

Tuy nhiên, các mô hình học sâu cũng đối mặt với nhiều thách thức: làm thế nào để xử lý dữ liệu mất cân bằng, làm thế nào để học được mối quan hệ giữa các đặc trưng giao dịch, và làm thế nào để khai thác thông tin tuần tự theo thời gian mà không tăng quá nhiều chi phí tính toán. Đó cũng chính là động lực để đồ án này đề xuất mô hình \textbf{MaRENet} – kết hợp nhiều kỹ thuật như \textbf{mask đặc trưng}, \textbf{phát triển giao dịch theo ngữ cảnh (transactional expansion)}, \textbf{mạng hồi tiếp (RNN)} và \textbf{hàm mất mát Focal Loss} nhằm tối ưu hiệu quả phát hiện gian lận trên dữ liệu thực tế.

\subsection{Những thách thức trong phát hiện gian lận}\label{subsec:fraud-pattern}
Phát hiện gian lận thẻ tín dụng là một bài toán phức tạp, không chỉ vì dữ liệu có tính chất phi cấu trúc và biến động, mà còn do yêu cầu khắt khe về tính thời gian thực và độ chính xác cao. Một hệ thống hiệu quả cần vượt qua nhiều rào cản về mặt kỹ thuật, dữ liệu và triển khai để có thể hoạt động ổn định trong môi trường tài chính thực tế, nơi các giao dịch diễn ra liên tục với khối lượng lớn.


\subsubsection{Sự mất cân bằng dữ liệu nghiêm trọng}
Trong các hệ thống giao dịch thực tế, tỷ lệ giao dịch gian lận thường rất nhỏ so với số lượng giao dịch hợp lệ, dẫn đến hiện tượng mất cân bằng lớp nghiêm trọng. Cụ thể, trong tập dữ liệu Credit Card Fraud từ Kaggle được sử dụng trong đồ án này, chỉ có \textbf{9.650} giao dịch gian lận trên tổng số \textbf{1.842.742} giao dịch. Điều này khiến các mô hình học máy dễ bị thiên lệch và bỏ sót các mẫu gian lận hiếm gặp. Hình~\ref{fig:fraud_ratio} minh họa tỷ lệ giao dịch hợp lệ và gian lận trong tập dữ liệu.

\begin{figure}[!ht]
    \centering
    \includegraphics[width=11cm]{Images/rate.png}
    \caption{\textbf{Tỷ lệ giao dịch hợp lệ và gian lận trong tập dữ liệu.} Giao dịch gian lận chiếm tỷ lệ rất nhỏ so với hợp lệ.}
    \label{fig:fraud_ratio}
\end{figure}

\textbf{Hệ quả của mất cân bằng lớp:}
\begin{itemize}
    \item Mô hình có xu hướng dự đoán tất cả là hợp lệ để đạt độ chính xác cao nhưng vô nghĩa.
    \item Các chỉ số như accuracy không phản ánh đúng chất lượng mô hình.
    \item Cần sử dụng các chiến lược như tái cân bằng dữ liệu, Focal Loss~\cite{lin2017focal}, hoặc kỹ thuật tăng cường như SMOTE.
\end{itemize}

\subsubsection{Tính tuần tự và ngữ cảnh của dữ liệu}
Giao dịch thẻ tín dụng thường xảy ra theo chuỗi hành vi, không độc lập. Các hành vi gian lận có thể được nhận diện thông qua sự thay đổi bất thường trong hành vi của người dùng theo thời gian. Việc khai thác mối quan hệ theo chuỗi là rất quan trọng để nhận diện những thay đổi bất thường này~\cite{jurgovsky2018sequence}. Các mô hình như RNN hoặc GRU được đánh giá là phù hợp với bài toán phát hiện theo chuỗi.

Ngoài ra, nhiều nghiên cứu chỉ ra rằng việc khai thác đặc trưng thời gian như thời điểm giao dịch (giờ/ngày) và khoảng cách giữa các giao dịch cũng giúp cải thiện chất lượng phát hiện~\cite{bahnsen2016feature, liu2020adaptive}.

\subsubsection{Dữ liệu phi cấu trúc và đa chiều}

Trong nhiều bài toán thực tế như phát hiện gian lận, dữ liệu đầu vào thường mang tính phi cấu trúc và đa chiều, bao gồm nhiều loại thuộc tính khác nhau như thời gian, danh mục, số liên tục, văn bản hoặc các đặc trưng hành vi. Các đặc trưng này đến từ các miền (domain) khác nhau, có phân phối không đồng nhất và mức độ đóng góp không giống nhau đối với mục tiêu dự đoán. Điều này khiến việc xử lý thủ công hoặc chuẩn hóa đồng nhất trở nên không hiệu quả.

Để giải quyết vấn đề đó, MaRENet sử dụng cơ chế \textit{channel-wise masking} nhằm tự động học trọng số cho từng đặc trưng đầu vào, cho phép mô hình tập trung vào các chiều thông tin có liên quan cao và loại bỏ nhiễu không cần thiết. Cơ chế này giúp nâng cao khả năng thích ứng của mô hình với dữ liệu có cấu trúc không đồng nhất và góp phần cải thiện hiệu suất dự đoán trên tập dữ liệu mất cân bằng.

\subsubsection{Yêu cầu xử lý thời gian thực}

Trong các hệ thống tài chính, phát hiện gian lận cần được thực hiện gần như tức thời để kịp thời ngăn chặn các giao dịch bất thường~\cite{fiore2019using}. Điều này đặt ra yêu cầu không chỉ về độ chính xác cao mà còn về độ trễ suy luận thấp trong triển khai thực tế. Trong thiết kế của MaRENet, cấu trúc mạng hồi tiếp được xây dựng theo hướng tinh gọn, với kích thước cửa sổ chuỗi giới hạn nhằm giảm thiểu chi phí tính toán.  

Tuy chưa có số liệu đo lường định lượng về độ trễ suy luận (inference latency), nhưng thiết kế tối giản này có tiềm năng phù hợp với các hệ thống yêu cầu xử lý nhanh. Đánh giá hiệu năng thời gian thực là một hướng mở cần được xem xét trong các nghiên cứu tiếp theo.

\subsubsection{Gian lận thay đổi liên tục theo thời gian}
Các hình thức gian lận luôn thay đổi để tránh bị phát hiện. Các mô hình tĩnh sẽ nhanh chóng lỗi thời nếu không được cập nhật. Việc học hành vi qua thời gian, như trong STGN~\cite{xie2024stgn}, giúp mô hình thích nghi với sự thay đổi về mặt thời gian và không gian trong hành vi giao dịch.

\subsubsection{Gian lận có tổ chức và mối liên kết ngầm}
Nhiều hành vi gian lận là có tổ chức, không ngẫu nhiên. Các nghiên cứu gần đây đã sử dụng mô hình học đồ thị để phát hiện mối liên hệ ngầm giữa người dùng, IP, thiết bị,… . Các phương pháp như GCN~\cite{dou2020enhancing}, GraphSAGE~\cite{weber2018scalable} hay DeepWalk~\cite{perozzi2014deepwalk} chứng minh hiệu quả trong việc phát hiện gian lận có tổ chức.

Tóm lại, bài toán phát hiện gian lận là một bài toán đầy thử thách, yêu cầu mô hình không chỉ chính xác mà còn thích nghi tốt, linh hoạt và có hiệu suất cao để ứng dụng thực tế.

\subsection{Động lực phát triển mô hình MaRENet}

Mặc dù các mô hình học sâu hiện đại đã đạt được nhiều tiến bộ trong việc phát hiện gian lận, nhưng việc áp dụng trực tiếp các kiến trúc phổ biến như MLP, CNN hay các mô hình hồi tiếp như RNN vào dữ liệu giao dịch thực tế vẫn còn tồn tại nhiều hạn chế. Cụ thể, các mô hình này thường được thiết kế để xử lý dữ liệu tuần tự hoặc hình ảnh tiêu chuẩn, trong khi dữ liệu giao dịch tài chính lại có cấu trúc không đồng nhất, phân bố theo thời gian và chứa nhiều đặc trưng phân loại có tính phi tuyến cao.

Thêm vào đó, như đã trình bày trong mục ~\ref{subsec:fraud-pattern}, các hành vi gian lận không diễn ra một cách độc lập mà thường nằm trong các chuỗi hành vi phức tạp. Do vậy, một mô hình phát hiện gian lận hiệu quả không chỉ cần xử lý được tính phi cấu trúc và tính tuần tự của dữ liệu, mà còn phải đảm bảo khả năng nhận diện những thay đổi hành vi tinh vi của người dùng trong ngữ cảnh giao dịch.

Trong những năm gần đây, một số kiến trúc tiên tiến như TabTransformer~\cite{huang2020tabtransformer}, STGNN~\cite{xie2024stgn}, hoặc STAN~\cite{cheng2020spatio} đã được áp dụng thành công trong bài toán phát hiện gian lận nhờ khả năng học đặc trưng mạnh mẽ và khả năng mô hình hóa mối quan hệ phi tuyến. Tuy nhiên, hầu hết các kiến trúc này đều yêu cầu tài nguyên tính toán lớn và khó triển khai trong môi trường thời gian thực. Các mô hình học đồ thị như GCN hoặc GAT cũng cần xây dựng biểu đồ liên kết phức tạp giữa các thực thể giao dịch, gây khó khăn trong xử lý trực tuyến. Điều này thúc đẩy nhu cầu về một mô hình vừa hiệu quả, vừa nhẹ và dễ triển khai – là tiền đề cho sự ra đời của MaRENet.

Trong bối cảnh đó, mô hình \textbf{MaRENet} được đề xuất với mục tiêu giải quyết trực tiếp các thách thức này thông qua việc tích hợp 4 thành phần cốt lõi:

\begin{itemize}
    \item \textbf{Transactional Expansion:} giúp mở rộng ngữ cảnh hành vi bằng cách kết hợp các giao dịch trước đó của người dùng thành một chuỗi có độ dài cố định, đảm bảo mô hình có đủ thông tin để học hành vi theo thời gian.
    \item \textbf{Channel-wise Feature Mask:} là cơ chế học trọng số cho từng đặc trưng đầu vào, giúp mô hình tập trung vào các đặc trưng quan trọng, đồng thời loại bỏ các đặc trưng nhiễu trong từng phiên giao dịch cụ thể.
    \item \textbf{Recurrent Neural Network (RNN):} giúp mô hình hóa mối quan hệ tuần tự trong chuỗi giao dịch, nhận diện các chuỗi hành vi bất thường theo thời gian.
    \item \textbf{Focal Loss:} là thành phần quan trọng giúp mô hình xử lý hiệu quả bài toán mất cân bằng lớp, nâng cao độ nhạy với các mẫu gian lận hiếm gặp.
\end{itemize}

Từng thành phần trong kiến trúc MaRENet đều đóng vai trò quan trọng đối với bài toán phát hiện gian lận thẻ tín dụng (CCFD). Cụ thể, cơ chế làm nổi bật đặc trưng giúp mô hình tập trung vào các yếu tố bất thường trong giao dịch; mạng hồi tiếp học được hành vi người dùng theo thời gian; và hàm mất mát phù hợp đảm bảo khả năng xử lý dữ liệu mất cân bằng, đặc trưng của các hệ thống tài chính thực tế.
Transactional Expansion cung cấp bối cảnh lịch sử; Feature Mask xử lý tính không đồng nhất của dữ liệu; RNN khai thác thông tin thời gian; còn Focal Loss giải quyết mất cân bằng nhãn. Sự phối hợp này đảm bảo rằng mô hình không chỉ học tốt mà còn có khả năng thích nghi với môi trường thực tế.

Hơn nữa, kiến trúc của MaRENet được thiết kế linh động, cho phép dễ dàng điều chỉnh để phù hợp với nhiều bài toán tài chính khác nhau, chẳng hạn như phát hiện rửa tiền (AML), theo dõi giao dịch chứng khoán bất thường, hoặc phân tích hành vi người dùng trong hệ thống ngân hàng số.

Một điểm đáng chú ý khác là tính giải thích của mô hình. Cơ chế Feature Mask cho phép phân tích tầm quan trọng của từng đặc trưng trong từng giao dịch, tạo điều kiện để chuyên gia điều tra hiểu được lý do phía sau mỗi cảnh báo. Điều này là yếu tố then chốt để xây dựng lòng tin khi triển khai các mô hình AI trong ngành tài chính.

Tóm lại, MaRENet ra đời từ nhu cầu thực tiễn trong việc phát hiện gian lận thẻ tín dụng trên dữ liệu lớn, không đồng đều, và luôn biến đổi. Động lực phát triển mô hình này không chỉ đến từ việc khắc phục các điểm yếu của các kiến trúc truyền thống, mà còn từ mục tiêu xây dựng một giải pháp khả thi, hiệu quả và phù hợp với triển khai thực tế tại các tổ chức tài chính và ngân hàng.

\subsection{Mục tiêu và phương pháp nghiên cứu}
\subsubsection{Mục tiêu của đồ án tốt nghiệp}

Trên cơ sở phân tích các thách thức trong bài toán phát hiện gian lận thẻ tín dụng và động lực xây dựng mô hình MaRENet, đồ án tốt nghiệp này hướng đến các mục tiêu cụ thể sau:

\begin{itemize}
    \item Nghiên cứu và phát triển mô hình MaRENet – một kiến trúc học sâu tích hợp cơ chế mở rộng giao dịch (Transactional Expansion), mặt nạ đặc trưng (Channel-wise Feature Mask), RNN và Focal Loss – nhằm cải thiện khả năng phát hiện các giao dịch gian lận hiếm gặp.
    \item Tối ưu hóa độ chính xác và đặc biệt là F1-score của mô hình, đảm bảo khả năng nhận diện chính xác các giao dịch bất thường trong điều kiện dữ liệu mất cân bằng nghiêm trọng.
    \item Thử nghiệm và đánh giá hiệu quả mô hình trên tập dữ liệu thực tế Credit Card Fraud từ Kaggle, với quy mô lớn (hơn 1.8 triệu giao dịch) và tỷ lệ gian lận rất thấp.
\end{itemize}

\subsubsection{Phương pháp đề xuất nghiên cứu}

Để hiện thực hóa các mục tiêu trên, đồ án triển khai một quy trình nghiên cứu theo hướng thực nghiệm kết hợp đánh giá định lượng. Cụ thể:

\begin{itemize}
    \item Triển khai một pipeline xử lý dữ liệu hoàn chỉnh bao gồm tiền xử lý, mở rộng ngữ cảnh giao dịch, huấn luyện và đánh giá mô hình.
    \item Áp dụng mô hình MaRENet trên tập huấn luyện, sau đó đánh giá bằng tập kiểm thử cố định để đảm bảo tính khách quan.
    \item Thực hiện so sánh với nhiều mô hình baseline, bao gồm các thuật toán học máy truyền thống (Random Forest, Logistic Regression) và các mô hình học sâu (LSTM, BiGRU), nhằm chứng minh hiệu quả tương đối của MaRENet.
    % \item Thiết kế mô hình hướng tới khả năng triển khai trong môi trường thời gian thực, thông qua tối ưu hóa chi phí tính toán và độ trễ xử lý.
\end{itemize}

Phương pháp nghiên cứu được thiết kế nhằm đánh giá mô hình không chỉ ở khía cạnh độ chính xác mà còn ở tính khả thi triển khai và tính ổn định khi ứng dụng vào thực tế.


\subsection{Tóm tắt chương 1}

Chương 1 đã trình bày tổng quan về bài toán phát hiện gian lận thẻ tín dụng trong bối cảnh thanh toán điện tử ngày càng phát triển, đồng thời làm rõ mức độ nghiêm trọng của các hành vi gian lận hiện đại với mức độ tinh vi và khó phát hiện ngày càng tăng. Các đặc thù phức tạp của dữ liệu giao dịch như tính phi cấu trúc, tính tuần tự, tính phi tuyến và sự mất cân bằng nghiêm trọng giữa số lượng giao dịch hợp lệ và gian lận đã được phân tích toàn diện, chỉ ra những thách thức cốt lõi trong việc thiết kế các hệ thống phát hiện gian lận hiệu quả.

Bên cạnh đó, chương cũng đã điểm qua những hạn chế của các phương pháp truyền thống và một số mô hình học sâu hiện tại khi được áp dụng trực tiếp vào dữ liệu giao dịch thực tế. Những bất cập này bao gồm độ trễ trong xử lý, thiếu khả năng thích ứng với các mẫu hành vi biến động theo thời gian và khó triển khai trong môi trường yêu cầu thời gian thực.

Từ thực trạng và các thách thức nêu trên, chương này đã giới thiệu động lực phát triển mô hình \textbf{MaRENet} như một hướng tiếp cận mới, với thiết kế mô đun hoá và tích hợp nhiều thành phần có tính chuyên biệt: \textbf{Transactional Expansion} giúp mở rộng ngữ cảnh hành vi, \textbf{Channel-wise Feature Masking} hỗ trợ tự động lựa chọn đặc trưng quan trọng, \textbf{RNN} học mối quan hệ tuần tự trong chuỗi hành vi người dùng, và \textbf{Focal Loss} tăng cường độ nhạy với các mẫu gian lận hiếm gặp.

Ngoài ra, chương 1 cũng đã xác lập rõ ràng mục tiêu nghiên cứu, phạm vi ứng dụng, và định hướng triển khai thực nghiệm của đồ án. Những nội dung này tạo nền tảng vững chắc để các chương tiếp theo đi sâu vào cơ sở lý thuyết liên quan (Chương 2), chi tiết kiến trúc và cơ chế hoạt động của mô hình đề xuất (Chương 3), và cuối cùng là quy trình đánh giá thực nghiệm và phân tích kết quả (Chương 4).